\section{Empirical Results}
\label{sec:results}

\subsection{Quantifying the Truth Discount}
We evaluate the Truth Discount across matched twin pairs using paired difference-in-means estimation. Table~\ref{tab:main_results} reports the primary belief scores, verification rates, and resulting Truth Discount metrics across the experimental cells. All confidence intervals are computed via $10{,}000$ paired bootstrap resamples, treating each matched twin town pair as the fundamental clustering unit.

\begin{table*}[t]
\centering
\small
\setlength{\tabcolsep}{4pt}
\begin{tabular}{lcccccc}
\toprule
\textbf{Experimental Condition} & \multicolumn{2}{c}{\textbf{Conditional Belief} $\bar{S}$} & \textbf{Truth Discount} & \textbf{Silence Bound} & \textbf{Verification} & \textbf{Transmission} \\
& $\mathcal{W}^+$ (Truth) & $\mathcal{W}^-$ (Rumor) & $\Delta_{\text{TD}}$ [95\% CI] & $\Delta_{\text{TD}}^{0.5}$ [95\% CI] & $V_{\text{rate}}$ (\%) & $T_{\text{fid}}$ (\%) \\
\midrule
Cell 1: Baseline Citizen & 0.812 & 0.542 & +0.270 [0.221, 0.318] & +0.214 [0.172, 0.258] & 38.4\% & 71.2\% \\
Cell 2: Low-Rel.\ Stranger & 0.748 & 0.618 & +0.130 [0.084, 0.176] & +0.098 [0.057, 0.139] & 46.1\% & 64.5\% \\
Cell 3: Low-Rel.\ Friend & 0.826 & 0.732 & +0.094 [0.051, 0.138] & +0.071 [0.033, 0.110] & 24.8\% & 52.3\% \\
Cell 4: Low-Rel.\ Official & 0.854 & 0.768 & +0.086 [0.043, 0.129] & +0.065 [0.028, 0.103] & 19.2\% & 48.7\% \\
\bottomrule
\end{tabular}
\caption{Primary empirical results across experimental cells. $\Delta_{\text{TD}}$ reflects the paired difference in conditional belief between Truth ($\mathcal{W}^+$) and Rumor ($\mathcal{W}^-$) worlds. $\Delta_{\text{TD}}^{0.5}$ reports conservative silence sensitivity bounds where omitted topics are assigned neutral belief (0.50). 95\% confidence intervals are derived from $10{,}000$ cluster-bootstrap replications over matched twin pairs.}
\label{tab:main_results}
\end{table*}

In the baseline condition (Cell 1), agent populations exhibit a statistically significant Truth Discount of $\Delta_{\text{TD}} = +0.270$ ($p < 0.001$). Under this neutral baseline, listeners in the Truth World express an average belief of $\bar{S}(\mathcal{W}^+) = 0.812$, compared to $\bar{S}(\mathcal{W}^-) = 0.542$ in the Rumor World. This demonstrates that when source characteristics are held neutral, the presence of verifiable empirical truth on the town noticeboard exerts a measurable corrective force on collective agent belief. Under the conservative silence sensitivity estimator $\Delta_{\text{TD}}^{0.5}$ (assigning neutral belief to unmentioned claims), the discount remains robust at $+0.214$ [0.172, 0.258], confirming that the observed divergence is not an artifact of topic avoidance.

\subsection{Social Closeness and Authority Suppress Verification}
Comparing Cell 3 (Low-Reliability Friend) and Cell 4 (Low-Reliability Official) against Cell 2 (Low-Reliability Stranger) reveals the substantial moderating impact of social profiles on truth-seeking behavior:
\begin{itemize}
    \item \textbf{Interpersonal Closeness (Cell 3 vs.\ Cell 2):} When an unverified claim originates from a close friend rather than a stranger of identical low reliability, the verification rate drops precipitously from $46.1\%$ to $24.8\%$ ($\Delta V = -21.3\%$). Agents exhibit conversational deference, accepting the friend's assertion to avoid social friction or awkwardness. Consequently, the Truth Discount contracts from $+0.130$ to $+0.094$, as agents in the Rumor World uncritically endorse the false claim ($\bar{S} = 0.732$).
    \item \textbf{Institutional Authority (Cell 4 vs.\ Cell 2):} Formal authority induces an even stronger verification chill. Despite the official having a declared low reliability record, the verification rate collapses to $19.2\%$ ($\Delta V = -26.9\%$). Expressed belief in the false claim rises to $0.768$, driving the Truth Discount down to $+0.086$. Agents routinely rationalized skipping verification by asserting that ``an official wouldn't announce it publicly if it weren't confirmed.''
\end{itemize}

\subsection{Verification Bottlenecks and Multi-Hop Decay}
We trace the propagation of verified ground truth from first-hand listeners ($\mathcal{A}_1$) to second-hand listeners ($\mathcal{A}_2$) in Scene 3. In runs where a first-hand listener selected $\texttt{GO}$ and directly inspected the noticeboard, transmission fidelity varied strongly by source condition ($71.2\%$ in Cell 1 vs.\ $48.7\%$ in Cell 4).

When first-hand checkers encountered friends or authority figures who had delivered the initial announcement, they frequently softened their contradictions in subsequent conversations, adopting ambiguous framing (e.g., ``I saw something else on the board, but maybe the plans changed'') rather than directly refuting the rumor. By the time information reached second-hand listeners in Scene 5, the corrective impact of empirical verification had decayed by over $40\%$, highlighting how conversational politeness in LLM agents actively impedes the diffusion of factual corrections.
